% Volume VII: Formal Synthesis
% Part XIII. A Mathematics of Disciplined Distrust
% Chapter 71: The Paranoia Functional
% Spine: proof-spines/ch071-spine.md

\chapter{The Paranoia Functional}
\label{ch:ch071}

This chapter introduces the \textbf{paranoia functional} as a formal measure of disciplined distrust: a belief-state should be evaluated not only by fit with appearances, but by its sensitivity to alternative causal histories that remain live under the evidence \(E\). The basic institutional measure is
\[
\mathcal P(E)
=
\sum_{H\in\mathcal H}
w(H)\mathbf 1[H\text{ remains admissible under }E].
\]
\Definition\ Here \(\mathcal H\) is the space of candidate histories, \(w(H)\) weights them, and the indicator records whether \(E\) has genuinely ruled a history out. A high value of \(\mathcal P(E)\) therefore need not mark confusion; it can mark justified refusal to collapse underdetermined cases too quickly.

A more refined version treats distrust through posterior dispersion rather than admissibility alone. The entropy-based variant is
\[
\mathcal P_H(E)
=
-\sum_{H\in\mathcal H}
P(H\mid E)\log P(H\mid E).
\]
\Definition\ This chapter's preface is thus simple: high entropy need not mean pathological uncertainty. Under evidentiary incompleteness, preserving structured multiplicity is part of epistemic responsibility, and \(\mathcal P(E)\) and \(\mathcal P_H(E)\) supply the formal objects on which the next chapters will define loss, reach, and bounded revision.
